\section*{Erd\H{o}s Problem \#274}

\subsection*{1) ROUND-2 OBJECTIVE}

Path (C), ending in a full proof of the original problem: prove the finite abelian version that appears in Erd\H{o}s's cited source, and separate it from the stronger modern webpage version for arbitrary groups.

This is the right direction because Round-1 already established only a special case (finite $p$-groups), while the original finite abelian question is in fact fully settled by a known theorem.

\subsection*{2) Round-1 FOUNDATION USED}

I use the following Round-1 facts.

\begin{enumerate}
\item Round-1 proved the statement for finite $p$-groups by a direct base-$p$ argument.
\item Round-1 identified the remaining gap correctly: the full arbitrary-group Herzog--Sch\"onheim conjecture was not proved there.
\end{enumerate}

\subsection*{3) NEW INSIGHT / TOOL (ROUND-2)}

The new ingredient is a precise source correction.

Erd\H{o}s explicitly wrote the question in the form:

\medskip
\centerline{\emph{``Let $G$ be a finite abelian group.  Can its elements be partitioned into cosets of distinct sizes?''}}
\medskip

Thus the original Erd\H{o}s version is \emph{finite abelian}, not arbitrary group.

The tool that closes this case is Sun's theorem on \emph{uniform covers by cosets of subnormal subgroups}:

\medskip
\centerline{\emph{In a nontrivial uniform cover of a group by cosets of subnormal subgroups, the finite indices cannot be pairwise distinct.}}
\medskip

Since every subgroup of an abelian group is normal, hence subnormal, this theorem applies immediately.

\subsection*{4) ATTACK PLAN (ROUND-2)}

The exact Round-1 gap was to pass from special cases to the full original finite abelian statement.

I will do that in three steps.

\begin{enumerate}
\item Restate the original finite abelian question precisely.
\item Invoke Sun's theorem in its exact form.
\item Verify all hypotheses: an exact partition is a nontrivial uniform cover of multiplicity $1$, and every subgroup in an abelian group is normal.
\end{enumerate}

\subsection*{5) WORK (ROUND-2)}

\paragraph{Original finite abelian statement.}
Let $G$ be a finite abelian group.  Suppose
\[
G=\bigcup_{i=1}^k a_iH_i
\qquad (k>1)
\]
is a disjoint union of left cosets of subgroups $H_i\le G$.

Write
\[
n_i=[G:H_i].
\]
Because $G$ is finite,
\[
|a_iH_i|=|H_i|=\frac{|G|}{n_i}.
\]
Therefore the coset sizes are pairwise distinct if and only if the indices $n_i$ are pairwise distinct.

So the problem is equivalent to asking whether a partition of a finite abelian group into cosets can have pairwise distinct indices.

\paragraph{Sun's theorem.}
Sun proved the following: if a group $G$ admits a \emph{nontrivial uniform cover}
\[
\{a_iG_i\}_{i=1}^k
\]
with finite indices
\[
[G:G_1]\le \cdots \le [G:G_k],
\]
and if all the subgroups $G_i$ are subnormal in $G$, then the indices cannot be pairwise distinct.

\paragraph{Verification of the hypotheses.}
Our partition
\[
G=\bigcup_{i=1}^k a_iH_i
\]
is a uniform cover of multiplicity $1$: every element of $G$ lies in exactly one listed coset.  Since $k>1$, it is nontrivial.

Because $G$ is abelian, every subgroup $H_i$ is normal in $G$, hence subnormal in $G$.

Also, $G$ is finite, so each index $[G:H_i]$ is finite.

Thus Sun's theorem applies directly and shows that the indices
\[
[G:H_1],\dots,[G:H_k]
\]
cannot be pairwise distinct.

By the size-index relation above, the coset sizes cannot be pairwise distinct either.

\paragraph{Conclusion.}
A finite abelian group cannot be partitioned into more than one coset of pairwise different sizes.

This proves the original Erd\H{o}s version completely.

\paragraph{What remains open.}
The stronger webpage version,
\[
\mbox{``If $G$ is an arbitrary group, can there exist an exact covering by cosets of different sizes?''}
\]
is the general Herzog--Sch\"onheim conjecture and remains open in full generality.

\subsection*{6) ADVERSARIAL VERIFICATION}

Potential weak points and their resolution:

\begin{enumerate}
\item \emph{Finite versus infinite groups.}  
The original Erd\H{o}s source is finite abelian, and the proof above uses finiteness to identify size with $|G|/[G:H_i]$.  So there is no hidden cardinal-arithmetic issue.
\item \emph{Does a partition count as a uniform cover?}  
Yes.  ``Uniform cover'' means every element is covered the same number of times; a partition is exactly the multiplicity-$1$ case.
\item \emph{Does subnormality hold?}  
Yes.  In an abelian group every subgroup is normal, hence automatically subnormal.
\item \emph{Did we accidentally prove the arbitrary-group statement?}  
No.  The final paragraph explicitly separates the solved finite abelian case from the still-open general-group version.
\end{enumerate}

\subsection*{7) FINAL}

\textbf{FULL SOLUTION --- FULL PROOF.}

The original Erd\H{o}s problem (finite abelian groups) is solved negatively: such a partition into cosets of distinct sizes does not exist.

\subsection*{8) COMPLETION ESTIMATE (MANDATORY)}

\textbf{COMPLETION: 100\%}

\subsection*{9) REFERENCES}

\begin{enumerate}
\item P. Erd\H{o}s, \emph{Some of my favourite problems in number theory}.  This source states the question for a finite abelian group.
\item T. F. Bloom, \emph{Erd\H{o}s Problem \#274}.  Current webpage statement and current status check.
\item Z.-W. Sun, \emph{On the Herzog--Sch\"onheim conjecture for uniform covers of groups}, J. Algebra 273 (2004), 153--175.
\item Round-1 output for Problem \#274 (complete proof for finite $p$-groups and gap identification for the general case).
\end{enumerate}



